% Options for packages loaded elsewhere
\PassOptionsToPackage{unicode}{hyperref}
\PassOptionsToPackage{hyphens}{url}
\documentclass[
]{article}
\usepackage{xcolor}
\usepackage[margin=1in]{geometry}
\usepackage{amsmath,amssymb}
\setcounter{secnumdepth}{-\maxdimen} % remove section numbering
\usepackage{iftex}
\ifPDFTeX
  \usepackage[T1]{fontenc}
  \usepackage[utf8]{inputenc}
  \usepackage{textcomp} % provide euro and other symbols
\else % if luatex or xetex
  \usepackage{unicode-math} % this also loads fontspec
  \defaultfontfeatures{Scale=MatchLowercase}
  \defaultfontfeatures[\rmfamily]{Ligatures=TeX,Scale=1}
\fi
\usepackage{lmodern}
\ifPDFTeX\else
  % xetex/luatex font selection
\fi
% Use upquote if available, for straight quotes in verbatim environments
\IfFileExists{upquote.sty}{\usepackage{upquote}}{}
\IfFileExists{microtype.sty}{% use microtype if available
  \usepackage[]{microtype}
  \UseMicrotypeSet[protrusion]{basicmath} % disable protrusion for tt fonts
}{}
\makeatletter
\@ifundefined{KOMAClassName}{% if non-KOMA class
  \IfFileExists{parskip.sty}{%
    \usepackage{parskip}
  }{% else
    \setlength{\parindent}{0pt}
    \setlength{\parskip}{6pt plus 2pt minus 1pt}}
}{% if KOMA class
  \KOMAoptions{parskip=half}}
\makeatother
\usepackage{color}
\usepackage{fancyvrb}
\newcommand{\VerbBar}{|}
\newcommand{\VERB}{\Verb[commandchars=\\\{\}]}
\DefineVerbatimEnvironment{Highlighting}{Verbatim}{commandchars=\\\{\}}
% Add ',fontsize=\small' for more characters per line
\usepackage{framed}
\definecolor{shadecolor}{RGB}{248,248,248}
\newenvironment{Shaded}{\begin{snugshade}}{\end{snugshade}}
\newcommand{\AlertTok}[1]{\textcolor[rgb]{0.94,0.16,0.16}{#1}}
\newcommand{\AnnotationTok}[1]{\textcolor[rgb]{0.56,0.35,0.01}{\textbf{\textit{#1}}}}
\newcommand{\AttributeTok}[1]{\textcolor[rgb]{0.13,0.29,0.53}{#1}}
\newcommand{\BaseNTok}[1]{\textcolor[rgb]{0.00,0.00,0.81}{#1}}
\newcommand{\BuiltInTok}[1]{#1}
\newcommand{\CharTok}[1]{\textcolor[rgb]{0.31,0.60,0.02}{#1}}
\newcommand{\CommentTok}[1]{\textcolor[rgb]{0.56,0.35,0.01}{\textit{#1}}}
\newcommand{\CommentVarTok}[1]{\textcolor[rgb]{0.56,0.35,0.01}{\textbf{\textit{#1}}}}
\newcommand{\ConstantTok}[1]{\textcolor[rgb]{0.56,0.35,0.01}{#1}}
\newcommand{\ControlFlowTok}[1]{\textcolor[rgb]{0.13,0.29,0.53}{\textbf{#1}}}
\newcommand{\DataTypeTok}[1]{\textcolor[rgb]{0.13,0.29,0.53}{#1}}
\newcommand{\DecValTok}[1]{\textcolor[rgb]{0.00,0.00,0.81}{#1}}
\newcommand{\DocumentationTok}[1]{\textcolor[rgb]{0.56,0.35,0.01}{\textbf{\textit{#1}}}}
\newcommand{\ErrorTok}[1]{\textcolor[rgb]{0.64,0.00,0.00}{\textbf{#1}}}
\newcommand{\ExtensionTok}[1]{#1}
\newcommand{\FloatTok}[1]{\textcolor[rgb]{0.00,0.00,0.81}{#1}}
\newcommand{\FunctionTok}[1]{\textcolor[rgb]{0.13,0.29,0.53}{\textbf{#1}}}
\newcommand{\ImportTok}[1]{#1}
\newcommand{\InformationTok}[1]{\textcolor[rgb]{0.56,0.35,0.01}{\textbf{\textit{#1}}}}
\newcommand{\KeywordTok}[1]{\textcolor[rgb]{0.13,0.29,0.53}{\textbf{#1}}}
\newcommand{\NormalTok}[1]{#1}
\newcommand{\OperatorTok}[1]{\textcolor[rgb]{0.81,0.36,0.00}{\textbf{#1}}}
\newcommand{\OtherTok}[1]{\textcolor[rgb]{0.56,0.35,0.01}{#1}}
\newcommand{\PreprocessorTok}[1]{\textcolor[rgb]{0.56,0.35,0.01}{\textit{#1}}}
\newcommand{\RegionMarkerTok}[1]{#1}
\newcommand{\SpecialCharTok}[1]{\textcolor[rgb]{0.81,0.36,0.00}{\textbf{#1}}}
\newcommand{\SpecialStringTok}[1]{\textcolor[rgb]{0.31,0.60,0.02}{#1}}
\newcommand{\StringTok}[1]{\textcolor[rgb]{0.31,0.60,0.02}{#1}}
\newcommand{\VariableTok}[1]{\textcolor[rgb]{0.00,0.00,0.00}{#1}}
\newcommand{\VerbatimStringTok}[1]{\textcolor[rgb]{0.31,0.60,0.02}{#1}}
\newcommand{\WarningTok}[1]{\textcolor[rgb]{0.56,0.35,0.01}{\textbf{\textit{#1}}}}
\usepackage{graphicx}
\makeatletter
\newsavebox\pandoc@box
\newcommand*\pandocbounded[1]{% scales image to fit in text height/width
  \sbox\pandoc@box{#1}%
  \Gscale@div\@tempa{\textheight}{\dimexpr\ht\pandoc@box+\dp\pandoc@box\relax}%
  \Gscale@div\@tempb{\linewidth}{\wd\pandoc@box}%
  \ifdim\@tempb\p@<\@tempa\p@\let\@tempa\@tempb\fi% select the smaller of both
  \ifdim\@tempa\p@<\p@\scalebox{\@tempa}{\usebox\pandoc@box}%
  \else\usebox{\pandoc@box}%
  \fi%
}
% Set default figure placement to htbp
\def\fps@figure{htbp}
\makeatother
\setlength{\emergencystretch}{3em} % prevent overfull lines
\providecommand{\tightlist}{%
  \setlength{\itemsep}{0pt}\setlength{\parskip}{0pt}}
\usepackage{bookmark}
\IfFileExists{xurl.sty}{\usepackage{xurl}}{} % add URL line breaks if available
\urlstyle{same}
% fallback for those not using the hyperref driver hyperxmp:
\makeatletter
\@ifundefined{xmpquote}{\newcommand{\xmpquote}[1]{#1}}{}
\makeatother
\hypersetup{
  pdftitle={HW01},
  pdfauthor={Jonathan Morris},
  hidelinks,
  pdfcreator={LaTeX via pandoc}}

\title{HW01}
\author{Jonathan Morris}
\date{2026-09-25}

\begin{document}
\maketitle

\subsection{Load Packages}\label{load-packages}

\begin{Shaded}
\begin{Highlighting}[]
\ImportTok{import}\NormalTok{ numpy }\ImportTok{as}\NormalTok{ np}
\ImportTok{import}\NormalTok{ pandas }\ImportTok{as}\NormalTok{ pd}
\ImportTok{import}\NormalTok{ matplotlib.pyplot }\ImportTok{as}\NormalTok{ plt}
\ImportTok{from}\NormalTok{ collections }\ImportTok{import}\NormalTok{ Counter}
\end{Highlighting}
\end{Shaded}

\begin{Shaded}
\begin{Highlighting}[]
\FunctionTok{library}\NormalTok{(dplyr)}
\end{Highlighting}
\end{Shaded}

\begin{verbatim}
## 
## Attaching package: 'dplyr'
\end{verbatim}

\begin{verbatim}
## The following objects are masked from 'package:stats':
## 
##     filter, lag
\end{verbatim}

\begin{verbatim}
## The following objects are masked from 'package:base':
## 
##     intersect, setdiff, setequal, union
\end{verbatim}

\begin{Shaded}
\begin{Highlighting}[]
\FunctionTok{library}\NormalTok{(tidyverse)}
\end{Highlighting}
\end{Shaded}

\begin{verbatim}
## -- Attaching core tidyverse packages ------------------------ tidyverse 2.0.0 --
## v forcats   1.0.0     v readr     2.1.5
## v ggplot2   3.5.2     v stringr   1.5.1
## v lubridate 1.9.4     v tibble    3.3.0
## v purrr     1.2.0     v tidyr     1.3.1
\end{verbatim}

\begin{verbatim}
## -- Conflicts ------------------------------------------ tidyverse_conflicts() --
## x dplyr::filter() masks stats::filter()
## x dplyr::lag()    masks stats::lag()
## i Use the conflicted package (<http://conflicted.r-lib.org/>) to force all conflicts to become errors
\end{verbatim}

\begin{Shaded}
\begin{Highlighting}[]
\FunctionTok{library}\NormalTok{(ggplot2)}
\end{Highlighting}
\end{Shaded}

\section{Question 7:}\label{question-7}

\subsection{(a) Read the CSV with ID as text. Preserve the severity
order Mild \textless{} Moderate \textless{} Severe. Print the dataset
dimensions, duplicate-ID count, missing-wait count, and negative-wait
count. (2
points)}\label{a-read-the-csv-with-id-as-text.-preserve-the-severity-order-mild-moderate-severe.-print-the-dataset-dimensions-duplicate-id-count-missing-wait-count-and-negative-wait-count.-2-points}

\begin{Shaded}
\begin{Highlighting}[]
\CommentTok{\#\# Read data:}
\NormalTok{d }\OperatorTok{=}\NormalTok{ pd.read\_csv(}\StringTok{"homework1\_clinic.csv"}\NormalTok{, dtype}\OperatorTok{=}\NormalTok{\{}\StringTok{"id"}\NormalTok{: }\StringTok{"string"}\NormalTok{\}) }

\CommentTok{\#\# Preserve severity order:}
\NormalTok{ordered\_levels }\OperatorTok{=}\NormalTok{ [}\StringTok{"Mild"}\NormalTok{, }\StringTok{"Moderate"}\NormalTok{, }\StringTok{"Severe"}\NormalTok{]}
\NormalTok{d[}\StringTok{"severity"}\NormalTok{] }\OperatorTok{=}\NormalTok{ pd.Categorical(d[}\StringTok{"severity"}\NormalTok{], categories }\OperatorTok{=}\NormalTok{ ordered\_levels, ordered }\OperatorTok{=} \VariableTok{True}\NormalTok{)}
\NormalTok{d\_sorted }\OperatorTok{=}\NormalTok{ d.sort\_values(by }\OperatorTok{=} \StringTok{"severity"}\NormalTok{)}
\BuiltInTok{print}\NormalTok{(d\_sorted)}
\end{Highlighting}
\end{Shaded}

\begin{verbatim}
##     id program  severity  age_years  prior_visits  wait_min  improved
## 0   01     Yes      Mild         23             0       2.0         1
## 6   07      No      Mild         25             0      10.0         0
## 7   08      No      Mild         37             1      10.0         1
## 1   02     Yes  Moderate         31             1       4.0         1
## 3   04     Yes  Moderate         29             2       6.0         1
## 8   09      No  Moderate         43             2      12.0         0
## 9   10      No  Moderate         66             5      25.0         0
## 2   03     Yes    Severe         45             3       4.0         0
## 4   05     Yes    Severe         60             4       8.0         0
## 5   06     Yes    Severe         52             2       NaN         1
## 10  11      No    Severe         48             2       NaN         1
## 11  12      No    Severe         57             3       NaN         0
\end{verbatim}

\begin{Shaded}
\begin{Highlighting}[]
\CommentTok{\#\# Print data shape:}
\BuiltInTok{print}\NormalTok{(d.shape) }
\end{Highlighting}
\end{Shaded}

\begin{verbatim}
## (12, 7)
\end{verbatim}

\begin{Shaded}
\begin{Highlighting}[]
\CommentTok{\#\# Print duplicated ID\textquotesingle{}s:}
\NormalTok{duplicated\_id }\OperatorTok{=}\NormalTok{ d[}\StringTok{\textquotesingle{}id\textquotesingle{}}\NormalTok{].duplicated().}\BuiltInTok{sum}\NormalTok{()}
\BuiltInTok{print}\NormalTok{(duplicated\_id)}
\end{Highlighting}
\end{Shaded}

\begin{verbatim}
## 0
\end{verbatim}

\begin{Shaded}
\begin{Highlighting}[]
\CommentTok{\#\# Print duplicated waits:}
\NormalTok{duplicated\_wait }\OperatorTok{=}\NormalTok{ d[}\StringTok{\textquotesingle{}wait\_min\textquotesingle{}}\NormalTok{].duplicated().}\BuiltInTok{sum}\NormalTok{()}
\BuiltInTok{print}\NormalTok{(duplicated\_wait)}
\end{Highlighting}
\end{Shaded}

\begin{verbatim}
## 4
\end{verbatim}

\begin{Shaded}
\begin{Highlighting}[]
\CommentTok{\#\# Print negative waits:}
\NormalTok{negative\_wait }\OperatorTok{=}\NormalTok{ (d[}\StringTok{\textquotesingle{}wait\_min\textquotesingle{}}\NormalTok{] }\OperatorTok{\textless{}} \DecValTok{0}\NormalTok{).}\BuiltInTok{sum}\NormalTok{()}
\BuiltInTok{print}\NormalTok{(negative\_wait)}
\end{Highlighting}
\end{Shaded}

\begin{verbatim}
## 0
\end{verbatim}

\subsection{(b) For observed waits, report n, mean, median, sample
variance and SD, Q1, Q3, IQR, fences, and flagged values. Use ddof=1 and
NumPy quantile method=``linear''. Compare the overlapping results with
your hand calculations in Question 4. (3
points)}\label{b-for-observed-waits-report-n-mean-median-sample-variance-and-sd-q1-q3-iqr-fences-and-flagged-values.-use-ddof1-and-numpy-quantile-methodlinear.-compare-the-overlapping-results-with-your-hand-calculations-in-question-4.-3-points}

\begin{Shaded}
\begin{Highlighting}[]
\CommentTok{\# Print n:}
\BuiltInTok{print}\NormalTok{(d[}\StringTok{"wait\_min"}\NormalTok{].size)}
\end{Highlighting}
\end{Shaded}

\begin{verbatim}
## 12
\end{verbatim}

\begin{Shaded}
\begin{Highlighting}[]
\CommentTok{\# Print the mean:}
\BuiltInTok{print}\NormalTok{(np.mean(d[}\StringTok{"wait\_min"}\NormalTok{]))}
\end{Highlighting}
\end{Shaded}

\begin{verbatim}
## 9.0
\end{verbatim}

\begin{Shaded}
\begin{Highlighting}[]
\CommentTok{\# Print the variance:}
\BuiltInTok{print}\NormalTok{(np.var(d[}\StringTok{"wait\_min"}\NormalTok{], ddof }\OperatorTok{=} \DecValTok{1}\NormalTok{))}
\end{Highlighting}
\end{Shaded}

\begin{verbatim}
## 47.0
\end{verbatim}

\begin{Shaded}
\begin{Highlighting}[]
\CommentTok{\# Print the variance:}
\BuiltInTok{print}\NormalTok{(np.std(d[}\StringTok{"wait\_min"}\NormalTok{], ddof }\OperatorTok{=} \DecValTok{1}\NormalTok{))}
\end{Highlighting}
\end{Shaded}

\begin{verbatim}
## 6.855654600401044
\end{verbatim}

\begin{Shaded}
\begin{Highlighting}[]
\CommentTok{\# Print the quartiles:}
\BuiltInTok{print}\NormalTok{(np.nanquantile(d[}\StringTok{"wait\_min"}\NormalTok{], [}\FloatTok{0.25}\NormalTok{, }\FloatTok{0.5}\NormalTok{, }\FloatTok{0.75}\NormalTok{], method }\OperatorTok{=} \StringTok{"linear"}\NormalTok{))}
\end{Highlighting}
\end{Shaded}

\begin{verbatim}
## [ 4.  8. 10.]
\end{verbatim}

\begin{Shaded}
\begin{Highlighting}[]
\CommentTok{\#\# Find the IQR and fences:}
\NormalTok{q1 }\OperatorTok{=}\NormalTok{ np.nanquantile(d[}\StringTok{"wait\_min"}\NormalTok{], [}\FloatTok{0.25}\NormalTok{], method }\OperatorTok{=} \StringTok{"linear"}\NormalTok{)}
\NormalTok{q3 }\OperatorTok{=}\NormalTok{ np.nanquantile(d[}\StringTok{"wait\_min"}\NormalTok{], [}\FloatTok{0.75}\NormalTok{], method }\OperatorTok{=} \StringTok{"linear"}\NormalTok{)}
\NormalTok{iqr }\OperatorTok{=}\NormalTok{ q3 }\OperatorTok{{-}}\NormalTok{ q1}
\NormalTok{fence1 }\OperatorTok{=}\NormalTok{ q1 }\OperatorTok{{-}} \FloatTok{1.5}\OperatorTok{*}\NormalTok{iqr}
\NormalTok{fence2 }\OperatorTok{=}\NormalTok{ q3 }\OperatorTok{+} \FloatTok{1.5}\OperatorTok{*}\NormalTok{iqr}

\CommentTok{\# Print the flagged values:}
\NormalTok{flagged }\OperatorTok{=}\NormalTok{ []}
\ControlFlowTok{for}\NormalTok{ i }\KeywordTok{in}\NormalTok{ d[}\StringTok{"wait\_min"}\NormalTok{]:}
  \ControlFlowTok{if}\NormalTok{ i }\OperatorTok{\textgreater{}}\NormalTok{ fence2 }\KeywordTok{or}\NormalTok{ i }\OperatorTok{\textless{}}\NormalTok{ fence1:}
\NormalTok{    flagged.append(i)}

\BuiltInTok{print}\NormalTok{(flagged)}
\end{Highlighting}
\end{Shaded}

\begin{verbatim}
## [25.0]
\end{verbatim}

The results are \ldots. compared to Q4!!

\subsection{(c) Produce a labeled histogram using edges 0, 5, 10, 15,
20, 25, 30, and print the bin counts. Use bins closed on the left and
open on the right, except the final bin includes 30. Also produce a
labeled horizontal boxplot with the 1.5 × IQR rule. (3
points)}\label{c-produce-a-labeled-histogram-using-edges-0-5-10-15-20-25-30-and-print-the-bin-counts.-use-bins-closed-on-the-left-and-open-on-the-right-except-the-final-bin-includes-30.-also-produce-a-labeled-horizontal-boxplot-with-the-1.5-iqr-rule.-3-points}

\subsubsection{Create Histogram:}\label{create-histogram}

\begin{Shaded}
\begin{Highlighting}[]
\NormalTok{custom\_bins }\OperatorTok{=}\NormalTok{ np.array([}\DecValTok{0}\NormalTok{, }\DecValTok{5}\NormalTok{, }\DecValTok{10}\NormalTok{, }\DecValTok{15}\NormalTok{, }\DecValTok{20}\NormalTok{, }\DecValTok{25}\NormalTok{, }\DecValTok{31}\NormalTok{])}

\NormalTok{plt.hist(d[}\StringTok{"wait\_min"}\NormalTok{], bins}\OperatorTok{=}\NormalTok{custom\_bins, edgecolor}\OperatorTok{=}\StringTok{\textquotesingle{}black\textquotesingle{}}\NormalTok{, color}\OperatorTok{=}\StringTok{\textquotesingle{}blue\textquotesingle{}}\NormalTok{)}

\CommentTok{\# Add tick labels:}
\NormalTok{centers }\OperatorTok{=}\NormalTok{ custom\_bins[:}\OperatorTok{{-}}\DecValTok{1}\NormalTok{] }\OperatorTok{+} \FloatTok{2.5} \CommentTok{\# Label bin centers except for the upper bound.}
\NormalTok{labels }\OperatorTok{=}\NormalTok{ [}\SpecialStringTok{f"}\SpecialCharTok{\{}\NormalTok{b}\SpecialCharTok{\}}\SpecialStringTok{{-}}\SpecialCharTok{\{}\NormalTok{b}\OperatorTok{+}\DecValTok{4}\SpecialCharTok{\}}\SpecialStringTok{"} \ControlFlowTok{for}\NormalTok{ b }\KeywordTok{in}\NormalTok{ custom\_bins[:}\OperatorTok{{-}}\DecValTok{1}\NormalTok{]] }\CommentTok{\# Create a list of bin labels.}
\NormalTok{labels[}\OperatorTok{{-}}\DecValTok{1}\NormalTok{] }\OperatorTok{=} \StringTok{"25{-}30"} \CommentTok{\# Add a last label for the final bin.}
\NormalTok{plt.xticks(centers, labels) }\CommentTok{\# Add the labels to the plot.}
\end{Highlighting}
\end{Shaded}

\begin{verbatim}
## ([<matplotlib.axis.XTick object at 0x110f59e80>, <matplotlib.axis.XTick object at 0x110f98910>, <matplotlib.axis.XTick object at 0x110f99090>, <matplotlib.axis.XTick object at 0x110f99810>, <matplotlib.axis.XTick object at 0x110f99f90>, <matplotlib.axis.XTick object at 0x110f9a710>], [Text(2.5, 0, '0-4'), Text(7.5, 0, '5-9'), Text(12.5, 0, '10-14'), Text(17.5, 0, '15-19'), Text(22.5, 0, '20-24'), Text(27.5, 0, '25-30')])
\end{verbatim}

\begin{Shaded}
\begin{Highlighting}[]
\NormalTok{plt.title(}\StringTok{"Histogram with Specific Bin Width (Size = 5)"}\NormalTok{)}
\NormalTok{plt.xlabel(}\StringTok{"Wait time (min)"}\NormalTok{)}
\NormalTok{plt.ylabel(}\StringTok{"Count"}\NormalTok{)}
\NormalTok{plt.show()}
\end{Highlighting}
\end{Shaded}

\pandocbounded{\includegraphics[keepaspectratio]{09242026_HW01_files/figure-latex/unnamed-chunk-5-1.pdf}}

\subsubsection{Create Boxplot:}\label{create-boxplot}

\begin{Shaded}
\begin{Highlighting}[]
\NormalTok{plt.boxplot(d[}\StringTok{"wait\_min"}\NormalTok{].dropna(), orientation}\OperatorTok{=}\StringTok{"horizontal"}\NormalTok{,}
\NormalTok{whis}\OperatorTok{=}\FloatTok{1.5}\NormalTok{, showmeans}\OperatorTok{=}\VariableTok{False}\NormalTok{)}
\end{Highlighting}
\end{Shaded}

\begin{verbatim}
## {'whiskers': [<matplotlib.lines.Line2D object at 0x111a07b10>, <matplotlib.lines.Line2D object at 0x111a07c50>], 'caps': [<matplotlib.lines.Line2D object at 0x111a07d90>, <matplotlib.lines.Line2D object at 0x111a07ed0>], 'boxes': [<matplotlib.lines.Line2D object at 0x111a079d0>], 'medians': [<matplotlib.lines.Line2D object at 0x111a24050>], 'fliers': [<matplotlib.lines.Line2D object at 0x111a24190>], 'means': []}
\end{verbatim}

\begin{Shaded}
\begin{Highlighting}[]
\NormalTok{plt.title(}\StringTok{"Boxplot Showing Distribution of Wait Times"}\NormalTok{)}
\NormalTok{plt.xlabel(}\StringTok{"Wait (minutes)"}\NormalTok{)}
\NormalTok{plt.show()}
\end{Highlighting}
\end{Shaded}

\pandocbounded{\includegraphics[keepaspectratio]{09242026_HW01_files/figure-latex/unnamed-chunk-6-3.pdf}}

\subsection{(d) In two or three sentences, describe the distribution,
explain why the mean uses the observed-wait count, and state how you
would handle the flagged wait. (2
points)}\label{d-in-two-or-three-sentences-describe-the-distribution-explain-why-the-mean-uses-the-observed-wait-count-and-state-how-you-would-handle-the-flagged-wait.-2-points}

The mean uses the observed wait count in the basic formula for sample
mean by dividing the sum of all the observations by the number of non-na
values (i.e., the observed wait count). In this case, since the flagged
wait is so much higher than other waits, I would run analyses with and
without the outlier to see if results changed. If I were presenting or
writing a paper, I would report both results.

\section{Question 8 An observational comparison in
R}\label{question-8-an-observational-comparison-in-r}

\subsection{a. Read the same CSV. Print the program-by-improved count
table. For each program group, report total participants, observed
improvement outcomes, number improved, improvement proportion, and
number of missing waits. Do not exclude an observed improvement outcome
because its waiting time is missing. (3
points)}\label{a.-read-the-same-csv.-print-the-program-by-improved-count-table.-for-each-program-group-report-total-participants-observed-improvement-outcomes-number-improved-improvement-proportion-and-number-of-missing-waits.-do-not-exclude-an-observed-improvement-outcome-because-its-waiting-time-is-missing.-3-points}

\begin{Shaded}
\begin{Highlighting}[]
\CommentTok{\# Load the data:}
\NormalTok{d }\OtherTok{\textless{}{-}} \FunctionTok{read\_csv}\NormalTok{(}\StringTok{"homework1\_clinic.csv"}\NormalTok{)}
\end{Highlighting}
\end{Shaded}

\begin{verbatim}
## Rows: 12 Columns: 7
## -- Column specification --------------------------------------------------------
## Delimiter: ","
## chr (3): id, program, severity
## dbl (4): age_years, prior_visits, wait_min, improved
## 
## i Use `spec()` to retrieve the full column specification for this data.
## i Specify the column types or set `show_col_types = FALSE` to quiet this message.
\end{verbatim}

\begin{Shaded}
\begin{Highlighting}[]
\CommentTok{\# Report total participants:}
\NormalTok{total\_pp }\OtherTok{=}\NormalTok{ d }\SpecialCharTok{|\textgreater{}} 
  \FunctionTok{filter}\NormalTok{(}\SpecialCharTok{!}\FunctionTok{is.na}\NormalTok{(id)) }\SpecialCharTok{|\textgreater{}} 
  \FunctionTok{count}\NormalTok{()}

\FunctionTok{paste0}\NormalTok{(}\StringTok{"We have "}\NormalTok{, total\_pp, }\StringTok{" unique participants."}\NormalTok{)}
\end{Highlighting}
\end{Shaded}

\begin{verbatim}
## [1] "We have 12 unique participants."
\end{verbatim}

\begin{Shaded}
\begin{Highlighting}[]
\CommentTok{\# Report observed improvement outcomes:}
\NormalTok{improved\_obs }\OtherTok{=}\NormalTok{ d }\SpecialCharTok{|\textgreater{}} 
  \FunctionTok{filter}\NormalTok{(}\SpecialCharTok{!}\FunctionTok{is.na}\NormalTok{(improved)) }\SpecialCharTok{|\textgreater{}} 
  \FunctionTok{count}\NormalTok{() }\SpecialCharTok{|\textgreater{}} 
  \FunctionTok{pull}\NormalTok{()}

\FunctionTok{paste0}\NormalTok{(}\StringTok{"We have "}\NormalTok{, improved\_obs, }\StringTok{" observed improvement outcomes."}\NormalTok{)}
\end{Highlighting}
\end{Shaded}

\begin{verbatim}
## [1] "We have 12 observed improvement outcomes."
\end{verbatim}

\begin{Shaded}
\begin{Highlighting}[]
\CommentTok{\# Report number improved:}
\NormalTok{improved\_count }\OtherTok{=}\NormalTok{ d }\SpecialCharTok{|\textgreater{}} 
  \FunctionTok{filter}\NormalTok{(improved }\SpecialCharTok{==} \DecValTok{1}\NormalTok{) }\SpecialCharTok{|\textgreater{}} 
  \FunctionTok{count}\NormalTok{() }\SpecialCharTok{|\textgreater{}} 
  \FunctionTok{pull}\NormalTok{()}

\FunctionTok{paste0}\NormalTok{(}\StringTok{"We have "}\NormalTok{, improved\_count, }\StringTok{" particpants that improved."}\NormalTok{)}
\end{Highlighting}
\end{Shaded}

\begin{verbatim}
## [1] "We have 6 particpants that improved."
\end{verbatim}

\begin{Shaded}
\begin{Highlighting}[]
\CommentTok{\# Report improvement proportion:}
\NormalTok{improved\_proportion }\OtherTok{=}\NormalTok{ improved\_count }\SpecialCharTok{/}\NormalTok{ improved\_obs }

\FunctionTok{paste0}\NormalTok{(}\StringTok{"Our improvement proportion is: "}\NormalTok{, improved\_proportion)}
\end{Highlighting}
\end{Shaded}

\begin{verbatim}
## [1] "Our improvement proportion is: 0.5"
\end{verbatim}

\begin{Shaded}
\begin{Highlighting}[]
\CommentTok{\# Report number of missing waits:}
\NormalTok{missing\_waits }\OtherTok{=}\NormalTok{ d }\SpecialCharTok{|\textgreater{}} 
  \FunctionTok{filter}\NormalTok{(}\FunctionTok{is.na}\NormalTok{(wait\_min)) }\SpecialCharTok{|\textgreater{}} 
  \FunctionTok{count}\NormalTok{() }\SpecialCharTok{|\textgreater{}} 
  \FunctionTok{pull}\NormalTok{()}

\FunctionTok{paste0}\NormalTok{(}\StringTok{"We have "}\NormalTok{, missing\_waits, }\StringTok{" missing wait times."}\NormalTok{)}
\end{Highlighting}
\end{Shaded}

\begin{verbatim}
## [1] "We have 3 missing wait times."
\end{verbatim}

\subsection{b. Calculate the Yes-minus-No difference in improvement
proportions. Express it as both a proportion difference and percentage
points. (2
points)}\label{b.-calculate-the-yes-minus-no-difference-in-improvement-proportions.-express-it-as-both-a-proportion-difference-and-percentage-points.-2-points}

\begin{Shaded}
\begin{Highlighting}[]
\CommentTok{\# Calculate yes count:}
\NormalTok{improved\_yes\_c }\OtherTok{=}\NormalTok{ d }\SpecialCharTok{|\textgreater{}} 
  \FunctionTok{filter}\NormalTok{(program }\SpecialCharTok{==} \StringTok{"Yes"}\NormalTok{, improved }\SpecialCharTok{==} \DecValTok{1}\NormalTok{) }\SpecialCharTok{|\textgreater{}} 
  \FunctionTok{count}\NormalTok{() }\SpecialCharTok{|\textgreater{}} 
  \FunctionTok{pull}\NormalTok{()}

\CommentTok{\# Calculate no count:}
\NormalTok{improved\_no\_c }\OtherTok{=}\NormalTok{ d }\SpecialCharTok{|\textgreater{}} 
  \FunctionTok{filter}\NormalTok{(program }\SpecialCharTok{==} \StringTok{"No"}\NormalTok{, improved }\SpecialCharTok{==} \DecValTok{1}\NormalTok{) }\SpecialCharTok{|\textgreater{}} 
  \FunctionTok{count}\NormalTok{() }\SpecialCharTok{|\textgreater{}} 
  \FunctionTok{pull}\NormalTok{()}

\CommentTok{\# Calculate ratios:}
\NormalTok{yes\_ratio }\OtherTok{=}\NormalTok{ improved\_yes\_c }\SpecialCharTok{/}\NormalTok{ d }\SpecialCharTok{|\textgreater{}} \FunctionTok{filter}\NormalTok{(program }\SpecialCharTok{==} \StringTok{"Yes"}\NormalTok{) }\SpecialCharTok{|\textgreater{}} \FunctionTok{count}\NormalTok{() }\SpecialCharTok{|\textgreater{}} \FunctionTok{pull}\NormalTok{()}
\NormalTok{no\_ratio }\OtherTok{=}\NormalTok{ improved\_no\_c }\SpecialCharTok{/}\NormalTok{ d }\SpecialCharTok{|\textgreater{}} \FunctionTok{filter}\NormalTok{(program }\SpecialCharTok{==} \StringTok{"No"}\NormalTok{) }\SpecialCharTok{|\textgreater{}} \FunctionTok{count}\NormalTok{() }\SpecialCharTok{|\textgreater{}} \FunctionTok{pull}\NormalTok{()}

\CommentTok{\# Calculate difference:}
\FunctionTok{paste0}\NormalTok{(}\StringTok{"The yes{-}no difference is: "}\NormalTok{, yes\_ratio }\SpecialCharTok{{-}}\NormalTok{ no\_ratio, }\StringTok{"."}\NormalTok{)}
\end{Highlighting}
\end{Shaded}

\begin{verbatim}
## [1] "The yes-no difference is: 0.333333333333333."
\end{verbatim}

\subsection{c.~Make a labeled bar chart of the two improvement
proportions with a vertical axis from 0 to 1. (1
point)}\label{c.-make-a-labeled-bar-chart-of-the-two-improvement-proportions-with-a-vertical-axis-from-0-to-1.-1-point}

\begin{Shaded}
\begin{Highlighting}[]
\CommentTok{\# Make a data frame with the proportions:}
\NormalTok{df }\OtherTok{=} \FunctionTok{tibble}\NormalTok{(}
  \AttributeTok{program =} \FunctionTok{c}\NormalTok{(}\StringTok{"Yes"}\NormalTok{, }\StringTok{"No"}\NormalTok{),}
  \AttributeTok{ratio =} \FunctionTok{c}\NormalTok{(yes\_ratio, no\_ratio)}
\NormalTok{)}

\NormalTok{df }\SpecialCharTok{|\textgreater{}} 
  \FunctionTok{ggplot}\NormalTok{(}\FunctionTok{aes}\NormalTok{(}\AttributeTok{x =}\NormalTok{ program, }\AttributeTok{y =}\NormalTok{ ratio)) }\SpecialCharTok{+}
  \FunctionTok{geom\_col}\NormalTok{() }\SpecialCharTok{+}
  \FunctionTok{ylim}\NormalTok{(}\DecValTok{0}\NormalTok{, }\DecValTok{1}\NormalTok{) }\SpecialCharTok{+} 
  \FunctionTok{xlab}\NormalTok{(}\StringTok{"Program"}\NormalTok{) }\SpecialCharTok{+}
  \FunctionTok{ylab}\NormalTok{(}\StringTok{"Improvement Proportion"}\NormalTok{) }\SpecialCharTok{+}
  \FunctionTok{ggtitle}\NormalTok{(}\StringTok{"Improvement Proportions by Program"}\NormalTok{)}
\end{Highlighting}
\end{Shaded}

\pandocbounded{\includegraphics[keepaspectratio]{09242026_HW01_files/figure-latex/unnamed-chunk-9-1.pdf}}

\subsection{d.~Write a defensible interpretation limited to these
participants. Explain why the difference is not automatically a causal
effect, naming one plausible baseline confounder and how it could affect
both program choice and improvement. (2
points)}\label{d.-write-a-defensible-interpretation-limited-to-these-participants.-explain-why-the-difference-is-not-automatically-a-causal-effect-naming-one-plausible-baseline-confounder-and-how-it-could-affect-both-program-choice-and-improvement.-2-points}

Based on these results, we can conclude that in this sample of 12
participants, those that chose counseling had almost double the
improvement ratio when compared to those students that did not chose
counseling. This is not necessarily a causal effect because participants
were not randomized and self-selected into groups. One possible
confounder is self-selection bias where participants who were more
likely to benefit from counseling could have selected into counseling
programs, therefore biasing the outcomes in number of participants that
improved.

\section{Question 9 Checking diagnostic probabilities in
Python}\label{question-9-checking-diagnostic-probabilities-in-python}

\section{In a hypothetical population, disease prevalence is 0.02,
sensitivity is 0.90, and specificity is 0.95. Let D mean disease and T
mean a positive test. Assume independent simulated people and the stated
conditional test
probabilities.}\label{in-a-hypothetical-population-disease-prevalence-is-0.02-sensitivity-is-0.90-and-specificity-is-0.95.-let-d-mean-disease-and-t-mean-a-positive-test.-assume-independent-simulated-people-and-the-stated-conditional-test-probabilities.}

\subsection{a. Use np.random.default\_rng(813109) and B=100000. First
generate B independent uniforms to assign disease; then generate B new
independent uniforms to assign positive tests with probability 0.90 for
people with disease and 0.05 otherwise. Print the four counts TP, FN,
FP, and TN, with clear labels. (3
points)}\label{a.-use-np.random.default_rng813109-and-b100000.-first-generate-b-independent-uniforms-to-assign-disease-then-generate-b-new-independent-uniforms-to-assign-positive-tests-with-probability-0.90-for-people-with-disease-and-0.05-otherwise.-print-the-four-counts-tp-fn-fp-and-tn-with-clear-labels.-3-points}

\begin{Shaded}
\begin{Highlighting}[]
\NormalTok{rng }\OperatorTok{=}\NormalTok{ np.random.default\_rng(}\DecValTok{813109}\NormalTok{)}
\NormalTok{B }\OperatorTok{=} \DecValTok{100000}

\CommentTok{\# Create disease vector}
\NormalTok{D }\OperatorTok{=}\NormalTok{ rng.random(B) }\OperatorTok{\textless{}} \FloatTok{0.02}

\CommentTok{\# Create test vector:}
\NormalTok{p\_pos }\OperatorTok{=}\NormalTok{ np.where(D, }\FloatTok{0.90}\NormalTok{, }\FloatTok{0.05}\NormalTok{) }\CommentTok{\# Creates an array which assigns 0.90 to any row with TRUE in D and 0.05 to FALSE, etc.}
\NormalTok{T }\OperatorTok{=}\NormalTok{ rng.random(B) }\OperatorTok{\textless{}}\NormalTok{ p\_pos}

\CommentTok{\# Calculate counts:}
\NormalTok{TP }\OperatorTok{=}\NormalTok{ np.}\BuiltInTok{sum}\NormalTok{(D }\OperatorTok{\&}\NormalTok{ T)}
\NormalTok{FN }\OperatorTok{=}\NormalTok{ np.}\BuiltInTok{sum}\NormalTok{(D }\OperatorTok{\&} \OperatorTok{\textasciitilde{}}\NormalTok{T)}
\NormalTok{FP }\OperatorTok{=}\NormalTok{ np.}\BuiltInTok{sum}\NormalTok{(}\OperatorTok{\textasciitilde{}}\NormalTok{D }\OperatorTok{\&}\NormalTok{ T)}
\NormalTok{TN }\OperatorTok{=}\NormalTok{ np.}\BuiltInTok{sum}\NormalTok{(}\OperatorTok{\textasciitilde{}}\NormalTok{D }\OperatorTok{\&} \OperatorTok{\textasciitilde{}}\NormalTok{T)}

\BuiltInTok{print}\NormalTok{(}\SpecialStringTok{f"True positives  (TP): }\SpecialCharTok{\{}\NormalTok{TP}\SpecialCharTok{\}}\SpecialStringTok{"}\NormalTok{)}
\end{Highlighting}
\end{Shaded}

\begin{verbatim}
## True positives  (TP): 1766
\end{verbatim}

\begin{Shaded}
\begin{Highlighting}[]
\BuiltInTok{print}\NormalTok{(}\SpecialStringTok{f"False negatives (FN): }\SpecialCharTok{\{}\NormalTok{FN}\SpecialCharTok{\}}\SpecialStringTok{"}\NormalTok{)}
\end{Highlighting}
\end{Shaded}

\begin{verbatim}
## False negatives (FN): 206
\end{verbatim}

\begin{Shaded}
\begin{Highlighting}[]
\BuiltInTok{print}\NormalTok{(}\SpecialStringTok{f"False positives (FP): }\SpecialCharTok{\{}\NormalTok{FP}\SpecialCharTok{\}}\SpecialStringTok{"}\NormalTok{)}
\end{Highlighting}
\end{Shaded}

\begin{verbatim}
## False positives (FP): 4913
\end{verbatim}

\begin{Shaded}
\begin{Highlighting}[]
\BuiltInTok{print}\NormalTok{(}\SpecialStringTok{f"True negatives  (TN): }\SpecialCharTok{\{}\NormalTok{TN}\SpecialCharTok{\}}\SpecialStringTok{"}\NormalTok{)}
\end{Highlighting}
\end{Shaded}

\begin{verbatim}
## True negatives  (TN): 93115
\end{verbatim}

\subsection{b. Calculate empirical sensitivity, specificity, PPV, and
NPV from their appropriate subsets. Compute the corresponding
theoretical probabilities using the given model and compare theory with
simulation. Briefly explain any differences; simulation is an empirical
check rather than a mathematical proof. (3
points)}\label{b.-calculate-empirical-sensitivity-specificity-ppv-and-npv-from-their-appropriate-subsets.-compute-the-corresponding-theoretical-probabilities-using-the-given-model-and-compare-theory-with-simulation.-briefly-explain-any-differences-simulation-is-an-empirical-check-rather-than-a-mathematical-proof.-3-points}

\begin{Shaded}
\begin{Highlighting}[]
\CommentTok{\# Calculate sensitivity, the proportion of people that test positive given that they have the disease:}
\NormalTok{sensitivity }\OperatorTok{=}\NormalTok{ TP}\OperatorTok{/}\NormalTok{(TP }\OperatorTok{+}\NormalTok{ FN)}
\BuiltInTok{print}\NormalTok{(}\StringTok{"The sensitivity is: "}\NormalTok{,sensitivity)}
\end{Highlighting}
\end{Shaded}

\begin{verbatim}
## The sensitivity is:  0.8955375253549696
\end{verbatim}

\begin{Shaded}
\begin{Highlighting}[]
\CommentTok{\# Calculate specificity, the proportion of people that test negative given that they do not have the disease:}
\NormalTok{specificity }\OperatorTok{=}\NormalTok{ TN}\OperatorTok{/}\NormalTok{(TN }\OperatorTok{+}\NormalTok{ FP)}
\BuiltInTok{print}\NormalTok{(}\StringTok{"The specificity is: "}\NormalTok{, specificity)}
\end{Highlighting}
\end{Shaded}

\begin{verbatim}
## The specificity is:  0.9498816664626434
\end{verbatim}

\begin{Shaded}
\begin{Highlighting}[]
\CommentTok{\# Calculate positive predictive value, the proportion of people that actually have the disease and test positive:}
\NormalTok{PPV }\OperatorTok{=}\NormalTok{ TP}\OperatorTok{/}\NormalTok{(TP }\OperatorTok{+}\NormalTok{ FP)}
\BuiltInTok{print}\NormalTok{(}\StringTok{"The PPV is: "}\NormalTok{, PPV)}
\end{Highlighting}
\end{Shaded}

\begin{verbatim}
## The PPV is:  0.2644108399460997
\end{verbatim}

\begin{Shaded}
\begin{Highlighting}[]
\CommentTok{\# Calculate negative predictive value, the proportion of people that test negative and do not have the disease:}
\NormalTok{NPV }\OperatorTok{=}\NormalTok{ TN}\OperatorTok{/}\NormalTok{(TN }\OperatorTok{+}\NormalTok{ FN)}
\BuiltInTok{print}\NormalTok{(}\StringTok{"The NPV is: "}\NormalTok{, NPV)}
\end{Highlighting}
\end{Shaded}

\begin{verbatim}
## The NPV is:  0.9977925654461482
\end{verbatim}

The results closely resemble what we would expect from a theoretical
calculation, with only slight differences in the computed sensitivity
and specificity values. These differences are largely accounted for when
we round to the second decimal place, but demonstrate that the
simulation will still have slight discrepancies when compared to the
theoretical case. As B goes to infinity, our theoretical and simulated
results will get closer and closer. This does not demonstrate a
mathematical proof because we only simulated the results in a heuristic
example. To prove these relationships, we would need to show that as B
goes to infinity, the simulated result converges to the theoretical
result.

\subsection{c.~Holding sensitivity and specificity fixed, calculate the
theoretical PPV when prevalence rises to 0.20. Explain the direction of
the change. (2
points)}\label{c.-holding-sensitivity-and-specificity-fixed-calculate-the-theoretical-ppv-when-prevalence-rises-to-0.20.-explain-the-direction-of-the-change.-2-points}

The PPV rises to 0.82 (see notes). This is likely due to the fact that
as the prevalence rises, we expect more people in the population to have
the disease, hence it is also much more likely that someone who tests
positive will have the disease also.

\section{Question 10 Sampling without replacement in
R}\label{question-10-sampling-without-replacement-in-r}

\section{A group contains eight distinct people labeled 1--8; people
1--3 have a specified allergy. In one trial, select three people
uniformly without replacement. The event of interest is that all three
selected people have the allergy. Reset the full eight-person population
before the next
trial.}\label{a-group-contains-eight-distinct-people-labeled-18-people-13-have-a-specified-allergy.-in-one-trial-select-three-people-uniformly-without-replacement.-the-event-of-interest-is-that-all-three-selected-people-have-the-allergy.-reset-the-full-eight-person-population-before-the-next-trial.}

\subsection{a. State the exact event probability as a product of
conditional probabilities. Also calculate the corresponding probability
for three independent draws with replacement. (2
points)}\label{a.-state-the-exact-event-probability-as-a-product-of-conditional-probabilities.-also-calculate-the-corresponding-probability-for-three-independent-draws-with-replacement.-2-points}

\subsection{b. Use B=100000 and set.seed(813110) to simulate groups with
sample.int(8, 3, replace=FALSE). Then use set.seed(813111) and
replace=TRUE for the second model. Report successes, empirical
probabilities, and theoretical probabilities for both models. (3
points)}\label{b.-use-b100000-and-set.seed813110-to-simulate-groups-with-sample.int8-3-replacefalse.-then-use-set.seed813111-and-replacetrue-for-the-second-model.-report-successes-empirical-probabilities-and-theoretical-probabilities-for-both-models.-3-points}

\begin{Shaded}
\begin{Highlighting}[]
\NormalTok{B }\OtherTok{=} \DecValTok{10000}
\FunctionTok{set.seed}\NormalTok{(}\DecValTok{813110}\NormalTok{)}
\CommentTok{\# Run the first model without replacement:}
\NormalTok{m1 }\OtherTok{\textless{}{-}} \FunctionTok{replicate}\NormalTok{(B, }\FunctionTok{all}\NormalTok{(}\FunctionTok{sample.int}\NormalTok{(}\DecValTok{8}\NormalTok{, }\DecValTok{3}\NormalTok{, }\AttributeTok{replace =} \ConstantTok{FALSE}\NormalTok{) }\SpecialCharTok{\textless{}=} \DecValTok{3}\NormalTok{))}
\NormalTok{succ1 }\OtherTok{\textless{}{-}} \FunctionTok{sum}\NormalTok{(m1)}
\NormalTok{emp1  }\OtherTok{\textless{}{-}} \FunctionTok{mean}\NormalTok{(m1)}
\NormalTok{theo1 }\OtherTok{\textless{}{-}}\NormalTok{ (}\DecValTok{3}\SpecialCharTok{/}\DecValTok{8}\NormalTok{) }\SpecialCharTok{*}\NormalTok{ (}\DecValTok{2}\SpecialCharTok{/}\DecValTok{7}\NormalTok{) }\SpecialCharTok{*}\NormalTok{ (}\DecValTok{1}\SpecialCharTok{/}\DecValTok{6}\NormalTok{)}

\FunctionTok{paste0}\NormalTok{(}\StringTok{"Successes for model 1: "}\NormalTok{, succ1)}
\end{Highlighting}
\end{Shaded}

\begin{verbatim}
## [1] "Successes for model 1: 161"
\end{verbatim}

\begin{Shaded}
\begin{Highlighting}[]
\FunctionTok{paste0}\NormalTok{(}\StringTok{"Empirical probabilitiy for model 1: "}\NormalTok{, emp1)}
\end{Highlighting}
\end{Shaded}

\begin{verbatim}
## [1] "Empirical probabilitiy for model 1: 0.0161"
\end{verbatim}

\begin{Shaded}
\begin{Highlighting}[]
\FunctionTok{paste0}\NormalTok{(}\StringTok{"Theoretical probability for model 1: "}\NormalTok{, theo1)}
\end{Highlighting}
\end{Shaded}

\begin{verbatim}
## [1] "Theoretical probability for model 1: 0.0178571428571429"
\end{verbatim}

\begin{Shaded}
\begin{Highlighting}[]
\FunctionTok{set.seed}\NormalTok{(}\DecValTok{813111}\NormalTok{)}
\CommentTok{\# Run the second model with replacement:}
\NormalTok{m2 }\OtherTok{\textless{}{-}} \FunctionTok{replicate}\NormalTok{(B, }\FunctionTok{all}\NormalTok{(}\FunctionTok{sample.int}\NormalTok{(}\DecValTok{8}\NormalTok{, }\DecValTok{3}\NormalTok{, }\AttributeTok{replace =} \ConstantTok{TRUE}\NormalTok{) }\SpecialCharTok{\textless{}=} \DecValTok{3}\NormalTok{))}
\NormalTok{succ2 }\OtherTok{\textless{}{-}} \FunctionTok{sum}\NormalTok{(m2)}
\NormalTok{emp2  }\OtherTok{\textless{}{-}} \FunctionTok{mean}\NormalTok{(m2)}
\NormalTok{theo2 }\OtherTok{\textless{}{-}}\NormalTok{ (}\DecValTok{3}\SpecialCharTok{/}\DecValTok{8}\NormalTok{)}\SpecialCharTok{\^{}}\DecValTok{3}
\FunctionTok{paste0}\NormalTok{(}\StringTok{"Successes for model 2: "}\NormalTok{, succ2)}
\end{Highlighting}
\end{Shaded}

\begin{verbatim}
## [1] "Successes for model 2: 527"
\end{verbatim}

\begin{Shaded}
\begin{Highlighting}[]
\FunctionTok{paste0}\NormalTok{(}\StringTok{"Empirical probabilitiy for model 2: "}\NormalTok{, emp2)}
\end{Highlighting}
\end{Shaded}

\begin{verbatim}
## [1] "Empirical probabilitiy for model 2: 0.0527"
\end{verbatim}

\begin{Shaded}
\begin{Highlighting}[]
\FunctionTok{paste0}\NormalTok{(}\StringTok{"Theoretical probability for model 2: "}\NormalTok{, theo2)}
\end{Highlighting}
\end{Shaded}

\begin{verbatim}
## [1] "Theoretical probability for model 2: 0.052734375"
\end{verbatim}

\subsection{c.~For the simulation without replacement, report running
empirical probabilities after 100, 1,000, 10,000, and 100,000 trials.
Plot the running probability against the number of trials, with a
horizontal line at the theoretical probability. A logarithmic horizontal
axis is acceptable. (3
points)}\label{c.-for-the-simulation-without-replacement-report-running-empirical-probabilities-after-100-1000-10000-and-100000-trials.-plot-the-running-probability-against-the-number-of-trials-with-a-horizontal-line-at-the-theoretical-probability.-a-logarithmic-horizontal-axis-is-acceptable.-3-points}

\begin{Shaded}
\begin{Highlighting}[]
\CommentTok{\# Set Seed:}
\FunctionTok{set.seed}\NormalTok{(}\DecValTok{813110}\NormalTok{)}

\CommentTok{\# Initialize data frames:}
\NormalTok{trials }\OtherTok{=} \FunctionTok{c}\NormalTok{(}\DecValTok{100}\NormalTok{, }\DecValTok{1000}\NormalTok{, }\DecValTok{10000}\NormalTok{, }\DecValTok{100000}\NormalTok{)}
\NormalTok{probs\_df }\OtherTok{=} \FunctionTok{data.frame}\NormalTok{(}\AttributeTok{B =} \FunctionTok{c}\NormalTok{(}\DecValTok{100}\NormalTok{, }\DecValTok{1000}\NormalTok{, }\DecValTok{10000}\NormalTok{, }\DecValTok{100000}\NormalTok{), }
                   \AttributeTok{probs =} \FunctionTok{c}\NormalTok{(}\ConstantTok{NA}\NormalTok{, }\ConstantTok{NA}\NormalTok{, }\ConstantTok{NA}\NormalTok{, }\ConstantTok{NA}\NormalTok{),}
                   \AttributeTok{log\_B =} \FunctionTok{log}\NormalTok{(B))}

\CommentTok{\# Run the simulations:}
\NormalTok{i }\OtherTok{=} \DecValTok{1}
\ControlFlowTok{for}\NormalTok{ (B }\ControlFlowTok{in}\NormalTok{ trials) \{}
    \CommentTok{\# Run Model}
\NormalTok{    m1 }\OtherTok{\textless{}{-}} \FunctionTok{replicate}\NormalTok{(B, }\FunctionTok{all}\NormalTok{(}\FunctionTok{sample.int}\NormalTok{(}\DecValTok{8}\NormalTok{, }\DecValTok{3}\NormalTok{, }\AttributeTok{replace =} \ConstantTok{FALSE}\NormalTok{) }\SpecialCharTok{\textless{}=} \DecValTok{3}\NormalTok{))}
\NormalTok{    succ1 }\OtherTok{\textless{}{-}} \FunctionTok{sum}\NormalTok{(m1)}
\NormalTok{    emp1  }\OtherTok{\textless{}{-}} \FunctionTok{mean}\NormalTok{(m1)}
    
    \CommentTok{\# Report results:}
    \FunctionTok{paste0}\NormalTok{(}\StringTok{"Successes for model 1 with B = "}\NormalTok{, B, }\StringTok{": "}\NormalTok{, succ1)}
    \FunctionTok{paste0}\NormalTok{(}\StringTok{"Empirical probabilitiy for model 1 with B = "}\NormalTok{, B, }\StringTok{": "}\NormalTok{, emp1)}
\NormalTok{    probs\_df[i, }\DecValTok{2}\NormalTok{] }\OtherTok{=}\NormalTok{ emp1}
\NormalTok{    i }\OtherTok{=}\NormalTok{ i }\SpecialCharTok{+} \DecValTok{1}
\NormalTok{\}}

\CommentTok{\# Plot the graph:}
\NormalTok{probs\_df }\SpecialCharTok{|\textgreater{}} 
  \FunctionTok{ggplot}\NormalTok{(}\FunctionTok{aes}\NormalTok{(}\AttributeTok{x =} \FunctionTok{log}\NormalTok{(B), }\AttributeTok{y =}\NormalTok{ probs)) }\SpecialCharTok{+}
  \FunctionTok{geom\_point}\NormalTok{() }\SpecialCharTok{+}
  \FunctionTok{geom\_line}\NormalTok{() }\SpecialCharTok{+}
  \FunctionTok{geom\_hline}\NormalTok{(}\AttributeTok{yintercept =}\NormalTok{ theo1, }\AttributeTok{linetype =} \StringTok{"dotted"}\NormalTok{) }\SpecialCharTok{+}
  \FunctionTok{xlab}\NormalTok{(}\StringTok{"Logarithm of Trial Numbers"}\NormalTok{) }\SpecialCharTok{+}
  \FunctionTok{ylab}\NormalTok{(}\StringTok{"Empirical Probabilities"}\NormalTok{) }\SpecialCharTok{+}
  \FunctionTok{annotate}\NormalTok{(}\StringTok{"text"}\NormalTok{, }\AttributeTok{x =} \DecValTok{6}\NormalTok{, }\AttributeTok{y =}\NormalTok{ theo1, }\AttributeTok{label =} \StringTok{"Theoretical Probability"}\NormalTok{,}
           \AttributeTok{hjust =} \DecValTok{1}\NormalTok{, }\AttributeTok{vjust =} \SpecialCharTok{{-}}\FloatTok{0.5}\NormalTok{, }\AttributeTok{color =} \StringTok{"red"}\NormalTok{)}
\end{Highlighting}
\end{Shaded}

\pandocbounded{\includegraphics[keepaspectratio]{09242026_HW01_files/figure-latex/unnamed-chunk-13-1.pdf}}

\subsection{d.~Explain why the two models have different probabilities.
Does the law of large numbers guarantee exact agreement after 25 trials,
or that each additional trial moves the estimate closer to the truth? (2
points)}\label{d.-explain-why-the-two-models-have-different-probabilities.-does-the-law-of-large-numbers-guarantee-exact-agreement-after-25-trials-or-that-each-additional-trial-moves-the-estimate-closer-to-the-truth-2-points}

The theoretical probability is defined for the case where our sample
size B trends to infinity. In theory, if we were to repeat this
experiment an infinite number of times, our empirical probability would
converge to the theoretical probability. This means that we will not
necessarily have exact agreement after 25 trials, and that as we keep
sampling greater and greater numbers, we get closer and closer to the
theoretical result.

\end{document}
